\section{Space Complexity --- Exp BFS --- Relax}
%%% generated by the blueprint pipeline from TCSlib.Complexity.SpaceComplexity.ExpBFS.Relax
\subsection*{Space Complexity --- Exp BFS --- Relax: Introduction}

Declarations of the Lean module \texttt{TCSlib.Complexity.SpaceComplexity.ExpBFS.Relax}.

\subsection*{Space Complexity --- Exp BFS --- Relax: Results}

\begin{definition}[Selected radius]
\lean{Complexity.LogProg.ExpBFS.rr}
\label{Complexity.LogProg.ExpBFS.rr}
\leanok
\uses{Complexity.LogProg.ExpBFS.St}
For a state $S$ and Boolean value $\varphi$, define the selected radius to be $S.s1$
when $\varphi$ is true and $S.s$ when $\varphi$ is false.
\end{definition}

\begin{lemma}[Selected Radius Component]
\lean{Complexity.LogProg.ExpBFS.toV_rad}
\label{Complexity.LogProg.ExpBFS.toV_rad}
\leanok
\uses{Complexity.LogProg.ExpBFS.St, Complexity.LogProg.ExpBFS.St.toV, Complexity.LogProg.ExpBFS.rad, Complexity.LogProg.ExpBFS.rr, Complexity.LogProg.ExpBFS.toPh}
\difficulty{1}
Let $S$ be an exponential BFS machine state and let $\phi$ be a Boolean. The entry of
the 17-component counter vector of $S$ at the index assigned to $\phi$ by the
phase-radius rule equals $S.s_1$ when $\phi$ is true and $S.s$ when $\phi$ is false.
\end{lemma}

\begin{definition}[Initial configuration predicate]
\lean{Complexity.LogProg.ExpBFS.ini}
\label{Complexity.LogProg.ExpBFS.ini}
\leanok
\uses{Turing.pairEncode}
Given an oracle
$\mathrm{orc}:\mathrm{Fin}(4)\to\mathrm{List}(\mathrm{Bool})\to\mathrm{Bool}$, a Boolean
input list $x$, and natural numbers $r,u$, define the initial predicate to be the
oracle's answer at index $1$ on the encoding of $x$ together with the binary
representations of $u$ and $r$, nested in that order.
\end{definition}

\begin{definition}[Step predicate]
\lean{Complexity.LogProg.ExpBFS.stp}
\label{Complexity.LogProg.ExpBFS.stp}
\leanok
\uses{Turing.pairEncode}
Given an oracle
$\mathrm{orc}:\mathrm{Fin}(4)\to\mathrm{List}(\mathrm{Bool})\to\mathrm{Bool}$, a Boolean
input list $x$, and natural numbers $r,u,v$, define the step predicate to be the
oracle's answer at index $2$ on the encoding of $x$ together with the binary
representations of $u$ and $v$, followed by the binary representation of $r$, nested in
that order.
\end{definition}

\begin{definition}[Cell value]
\lean{Complexity.LogProg.ExpBFS.cell}
\label{Complexity.LogProg.ExpBFS.cell}
\leanok
For natural numbers $b$ and $v$, define the cell value $c(b,v)$ to be the integer $b+v$.
\end{definition}

\begin{definition}[Monotonicity of marked positions]
\lean{Complexity.LogProg.ExpBFS.Mono}
\label{Complexity.LogProg.ExpBFS.Mono}
\leanok
For two Boolean-valued functions $M,M':\mathbb{Z}\to\{\mathrm{true},\mathrm{false}\}$,
say that $M'$ extends the marked positions of $M$ if every integer marked true by $M$ is
also marked true by $M'$.
\end{definition}

\begin{definition}[Agreement outside a memory region]
\lean{Complexity.LogProg.ExpBFS.Out}
\label{Complexity.LogProg.ExpBFS.Out}
\leanok
\uses{Complexity.LogProg.ExpBFS.cell}
For natural numbers $b,K$ and Boolean-valued functions
$M,M':\mathbb{Z}\to\{\mathrm{true},\mathrm{false}\}$, define $M'$ to agree with $M$
outside the $K$ positions with addresses $b+v$ for $0\le v<K$ if their values agree at
every integer $z$ that is not one of those addresses.
\end{definition}

\begin{lemma}[Reflexivity of Boolean implication]
\lean{Complexity.LogProg.ExpBFS.Mono.refl}
\label{Complexity.LogProg.ExpBFS.Mono.refl}
\leanok
\uses{Complexity.LogProg.ExpBFS.Mono}
\difficulty{1}
For every pair of integers indexed by a Boolean-valued function $M$, each value of $M$
that is true is true in $M$ itself.
\end{lemma}

\begin{lemma}[Transitivity of Boolean implication]
\lean{Complexity.LogProg.ExpBFS.Mono.trans}
\label{Complexity.LogProg.ExpBFS.Mono.trans}
\leanok
\uses{Complexity.LogProg.ExpBFS.Mono}
\difficulty{1}
Let $M$, $M'$, and $M''$ be Boolean-valued functions on the integers. If every integer
at which $M$ is true is also one at which $M'$ is true, and every integer at which $M'$
is true is also one at which $M''$ is true, then every integer at which $M$ is true is
one at which $M''$ is true.
\end{lemma}

\begin{lemma}[Reflexivity outside a finite interval]
\lean{Complexity.LogProg.ExpBFS.Out.refl}
\label{Complexity.LogProg.ExpBFS.Out.refl}
\leanok
\uses{Complexity.LogProg.ExpBFS.Out}
\difficulty{1}
Let $b,K$ be natural numbers, and let $M$ be a Boolean-valued function on the integers.
The function $M$ agrees with itself at every integer that is not equal to $b+v$ for any
natural number $v<K$.
\end{lemma}

\begin{lemma}[Transitivity outside a finite interval]
\lean{Complexity.LogProg.ExpBFS.Out.trans}
\label{Complexity.LogProg.ExpBFS.Out.trans}
\leanok
\uses{Complexity.LogProg.ExpBFS.Out}
\difficulty{1}
Let $b,K$ be natural numbers and let $M,M',M''$ be Boolean-valued functions on the
integers. Suppose $M$ and $M'$ agree at every integer other than the integers $b+v$ with
$v<K$, and $M'$ and $M''$ also agree at every such integer. Then $M$ and $M''$ agree at
every integer other than those integers.
\end{lemma}

\begin{lemma}[Monotonicity under setting a bit]
\lean{Complexity.LogProg.ExpBFS.mono_update}
\label{Complexity.LogProg.ExpBFS.mono_update}
\leanok
\uses{Complexity.LogProg.ExpBFS.Mono}
\difficulty{4}
For any map $M : \mathbb{Z} \to \{\mathrm{false},\mathrm{true}\}$ and integer $z$, every
integer at which $M$ is true is also one at which the map obtained by setting its value
at $z$ to true is true.
\end{lemma}

\begin{lemma}[Memory Update Outside the Marked Region]
\lean{Complexity.LogProg.ExpBFS.out_update}
\label{Complexity.LogProg.ExpBFS.out_update}
\leanok
\uses{Complexity.LogProg.ExpBFS.Out, Complexity.LogProg.ExpBFS.cell}
\difficulty{1}
Let $b,K,v\in\mathbb{N}$ and let $M:\mathbb{Z}\to\{\mathrm{false},\mathrm{true}\}$ be
memory contents. If $v<K$, then changing the memory value at cell $b+v$ to true leaves
every cell other than $b+u$ for $u<K$ unchanged.
\end{lemma}

\begin{lemma}[Injectivity of shifted cell indices]
\lean{Complexity.LogProg.ExpBFS.cell_inj}
\label{Complexity.LogProg.ExpBFS.cell_inj}
\leanok
\uses{Complexity.LogProg.ExpBFS.cell}
\difficulty{3}
For natural numbers $b,v,v'$, if the integers represented by $b+v$ and $b+v'$ are equal,
then $v=v'$.
\end{lemma}

\begin{lemma}[Concatenation of Guarded Runs]
\lean{Complexity.LogProg.ExpBFS.ex1_pre}
\label{Complexity.LogProg.ExpBFS.ex1_pre}
\leanok
\uses{Complexity.LogProg.ExpBFS.A, Complexity.LogProg.ExpBFS.G, Complexity.LogProg.ExpBFS.Lb, Complexity.LogProg.MConf, Complexity.LogProg.MReach, Complexity.LogProg.MReach.trans}
\difficulty{4}
Fix a memory machine, oracle, input, and guard predicate. If a configuration $a$ reaches
$b$ in finitely many steps with the guard holding at every time strictly before arrival,
and there exists $y$ such that $b$ reaches $f(y)$ with the guard holding at every time
strictly before arrival and $Q(y)$, then there exists $y$ such that $a$ reaches $f(y)$
with the guard holding at every time strictly before arrival and $Q(y)$.
\end{lemma}

\begin{lemma}[Concatenation of Guarded Runs with Two Outputs]
\lean{Complexity.LogProg.ExpBFS.ex2_pre}
\label{Complexity.LogProg.ExpBFS.ex2_pre}
\leanok
\uses{Complexity.LogProg.ExpBFS.A, Complexity.LogProg.ExpBFS.G, Complexity.LogProg.ExpBFS.Lb, Complexity.LogProg.MConf, Complexity.LogProg.MReach, Complexity.LogProg.MReach.trans}
\difficulty{4}
Fix a memory machine, oracle, input, and guard predicate. If a configuration $a$ reaches
$b$ in finitely many steps with the guard holding at every time strictly before arrival,
and there exist $y,z$ such that $b$ reaches $f(y,z)$ with the guard holding at every
time strictly before arrival and $Q(y,z)$, then there exist $y,z$ such that $a$ reaches
$f(y,z)$ with the guard holding at every time strictly before arrival and $Q(y,z)$.
\end{lemma}

\begin{definition}[Bounded reachability closure]
\lean{Complexity.LogProg.ExpBFS.Cl}
\label{Complexity.LogProg.ExpBFS.Cl}
\leanok
\uses{Complexity.LogProg.ExpBFS.ini, Complexity.LogProg.ExpBFS.stp}
Given an oracle
$\mathrm{orc}:\mathrm{Fin}(4)\to\mathrm{List}(\mathrm{Bool})\to\mathrm{Bool}$, a Boolean
input list $x$, and natural numbers $r,K,v$, define the bounded reachability predicate
to hold when there is some $u<K$ such that the initial predicate at $u$ is true and $v$
is reachable from $u$ by finitely many steps, each step satisfying the step predicate
and ending at a value below $K$.
\end{definition}

\begin{lemma}[Bound on reachable vertices]
\lean{Complexity.LogProg.ExpBFS.Cl.lt}
\label{Complexity.LogProg.ExpBFS.Cl.lt}
\leanok
\uses{Complexity.LogProg.ExpBFS.Cl}
\difficulty{4}
Let $r,K,v$ be natural numbers. If there is a vertex $u<K$ accepted by the
initial-vertex oracle query for $r$, and there is a finite sequence of zero or more
transitions from $u$ to $v$ in which each transition is accepted by the step oracle
query for $r$ and ends at a vertex less than $K$, then $v<K$.
\end{lemma}

\begin{lemma}[Closure under one transition]
\lean{Complexity.LogProg.ExpBFS.Cl.step}
\label{Complexity.LogProg.ExpBFS.Cl.step}
\leanok
\uses{Complexity.LogProg.ExpBFS.Cl, Complexity.LogProg.ExpBFS.stp}
\difficulty{3}
Let $r,K,u,v$ be natural numbers. If $u$ is reachable from some vertex below $K$
accepted by the initial-vertex oracle query, using a finite sequence of oracle-accepted
transitions whose endpoints are below $K$, and the step oracle accepts a transition from
$u$ to $v$ with $v<K$, then $v$ is reachable in the same sense.
\end{lemma}

\begin{definition}[State bounds]
\lean{Complexity.LogProg.ExpBFS.SBd}
\label{Complexity.LogProg.ExpBFS.SBd}
\leanok
\uses{Complexity.LogProg.ExpBFS.St, Complexity.LogProg.ExpBFS.St.Bd}
For natural numbers $B,P$ and a state $S$, define the state bounds condition to mean
that $S$ satisfies its $B$-bound, that $S.Bs+S.K\le B$ and $S.Bs+S.K\le P$, and that
$S.K+1\le B$.
\end{definition}

\begin{theorem}[Initial marking search]
\lean{Complexity.LogProg.ExpBFS.ini_run}
\label{Complexity.LogProg.ExpBFS.ini_run}
\leanok
\uses{Complexity.LogProg.ExpBFS.A, Complexity.LogProg.ExpBFS.G, Complexity.LogProg.ExpBFS.Mono, Complexity.LogProg.ExpBFS.Mono.refl, Complexity.LogProg.ExpBFS.Mono.trans, Complexity.LogProg.ExpBFS.Out, Complexity.LogProg.ExpBFS.Out.refl, Complexity.LogProg.ExpBFS.Out.trans, Complexity.LogProg.ExpBFS.SBd, Complexity.LogProg.ExpBFS.Site.idx, Complexity.LogProg.ExpBFS.Site.set, Complexity.LogProg.ExpBFS.St, Complexity.LogProg.ExpBFS.St.Bd, Complexity.LogProg.ExpBFS.St.toV_U, Complexity.LogProg.ExpBFS.acExit, Complexity.LogProg.ExpBFS.ac_run, Complexity.LogProg.ExpBFS.cell, Complexity.LogProg.ExpBFS.cell_inj, Complexity.LogProg.ExpBFS.cf, Complexity.LogProg.ExpBFS.ex1_pre, Complexity.LogProg.ExpBFS.ini, Complexity.LogProg.ExpBFS.mono_update, Complexity.LogProg.ExpBFS.out_update, Complexity.LogProg.ExpBFS.reach_cf, Complexity.LogProg.ExpBFS.rr, Complexity.LogProg.ExpBFS.step_of, Complexity.LogProg.ExpBFS.toV_rad, Complexity.LogProg.MReach, Complexity.LogProg.MReach.refl, Complexity.LogProg.MReach.trans, Complexity.LogProg.astep, Complexity.LogProg.vword_whole₂, Turing.pairEncode}
\difficulty{6}
Let $K_0,B,P\in\mathbb N$, let
$\mathrm{orc}:\{0,1,2,3\}\times\{0,1\}^*\to\{\mathrm{false},\mathrm{true}\}$ be an
oracle, let $x\in\{0,1\}^*$, and choose $\varphi\in\{\mathrm{false},\mathrm{true}\}$.
Let $j\in\mathbb N$, let $S=(N,L_g,s,s_1,X,w,K,B_s,U,V,T,T_2,Z,C,ch,c,c_1)$ be a state,
and let $M:\mathbb Z\to\{\mathrm{false},\mathrm{true}\}$. Suppose $K-U=j$, $U\le K$,
$T=T_2=Z=0$, every coordinate of $S$ is at most $B$, $B_s+K\le B$, $B_s+K\le P$, and
$K+1\le B$. Then there are a memory $M':\mathbb Z\to\{\mathrm{false},\mathrm{true}\}$
and a finite run of the specified machine from label $iU(\varphi)$, state $S$, memory
$M$, and memory-head position $0$ to label $rl0(\varphi)$, state obtained from $S$ by
replacing $U$ with $K$, memory $M'$, and head position $0$, such that every
configuration before the endpoint has all register values at most $B$, has memory-head
position between $0$ and $P$ inclusive, and satisfies the machine's instruction
preconditions. The memory $M'$ extends the marked positions of $M$, agrees with $M$
outside the cells $c(B_s,v)$ for $v<K$, and, for every $v$ with $U\le v<K$, marks
$c(B_s,v)$ whenever the oracle accepts the encoded input at radius $s_1$ if $\varphi$ is
true and at radius $s$ otherwise. Conversely, for every $v<K$, if $M'$ marks $c(B_s,v)$,
then either $M$ already marked that cell or the oracle accepts that encoded input at the
chosen radius.
\end{theorem}

\begin{theorem}[Inner relaxation run]
\lean{Complexity.LogProg.ExpBFS.inner_run}
\label{Complexity.LogProg.ExpBFS.inner_run}
\leanok
\uses{Complexity.LogProg.ExpBFS.A, Complexity.LogProg.ExpBFS.G, Complexity.LogProg.ExpBFS.Mono, Complexity.LogProg.ExpBFS.Mono.refl, Complexity.LogProg.ExpBFS.Mono.trans, Complexity.LogProg.ExpBFS.Out, Complexity.LogProg.ExpBFS.Out.refl, Complexity.LogProg.ExpBFS.Out.trans, Complexity.LogProg.ExpBFS.SBd, Complexity.LogProg.ExpBFS.Site.idx, Complexity.LogProg.ExpBFS.Site.set, Complexity.LogProg.ExpBFS.St, Complexity.LogProg.ExpBFS.St.Bd, Complexity.LogProg.ExpBFS.St.toV_V, Complexity.LogProg.ExpBFS.acExit, Complexity.LogProg.ExpBFS.ac_run, Complexity.LogProg.ExpBFS.cell, Complexity.LogProg.ExpBFS.cell_inj, Complexity.LogProg.ExpBFS.cf, Complexity.LogProg.ExpBFS.cont, Complexity.LogProg.ExpBFS.ex2_pre, Complexity.LogProg.ExpBFS.mono_update, Complexity.LogProg.ExpBFS.out_update, Complexity.LogProg.ExpBFS.reach_cf, Complexity.LogProg.ExpBFS.rr, Complexity.LogProg.ExpBFS.step_of, Complexity.LogProg.ExpBFS.stp, Complexity.LogProg.ExpBFS.toV_rad, Complexity.LogProg.MReach, Complexity.LogProg.MReach.refl, Complexity.LogProg.astep, Complexity.LogProg.vword_whole₃, Turing.pairEncode}
\difficulty{7}
Let $K_0,B,P$ be natural numbers, let
$\mathrm{orc}:\{0,1,2,3\}\times\{0,1\}^*\to\{\mathrm{false},\mathrm{true}\}$ be an
oracle, let $x$ be a finite Boolean string, and let $\varphi$ be a Boolean. For every
$u,j\in\mathbb N$, state $S=(N,L_g,s,s_1,X,w,K,B_s,U,V,T,T_2,Z,C,ch,c,c_1)$, and memory
$M:\mathbb Z\to\{\mathrm{false},\mathrm{true}\}$, suppose $K-V=j$, $V\le K$, $U=u<K$,
$T=T_2=Z=0$, $ch\le1$, every coordinate of $S$ is at most $B$, and $B_s+K\le B$,
$B_s+K\le P$, and $K+1\le B$. Then there are a memory $M'$ and $ch'\in\mathbb N$ such
that the machine $A(K_0)$, on input $x$ with oracle $\mathrm{orc}$, runs from state $S$,
memory $M$, head position $0$, and label corresponding to reading the mark of $V$ in the
relaxation round for $\varphi$, to the state obtained from $S$ by setting $V=0$ and
$ch=ch'$, memory $M'$, head position $0$, and label corresponding to incrementing $U$ in
that round. Every configuration before the endpoint has all register values at most $B$
and memory head position between $0$ and $P$. Moreover, $M'$ extends the marked
positions of $M$ and agrees with $M$ outside the cells $B_s+v$ for $v<K$, while
$ch'\le1$. For every $v<K$, if $M'$ marks $B_s+v$, then either $M$ already marks that
cell or $\mathrm{orc}$ accepts the encoded query for $(x,\rho,u,v)$, where $\rho=s_1$ if
$\varphi$ is true and $\rho=s$ otherwise; conversely, if $V\le v<K$ and that query is
accepted, then $M'$ marks $B_s+v$. If $M'=M$, then $ch'=ch$; if $M'\ne M$, then $ch'=1$.
\end{theorem}

\begin{lemma}[Antisymmetry of monotone maps]
\lean{Complexity.LogProg.ExpBFS.mono_antisymm}
\label{Complexity.LogProg.ExpBFS.mono_antisymm}
\leanok
\uses{Complexity.LogProg.ExpBFS.Mono}
\difficulty{3}
For maps $M,M' : \mathbb{Z} \to \{\mathrm{false},\mathrm{true}\}$, suppose every integer
at which $M$ is true is also one at which $M'$ is true, and conversely. Then $M=M'$.
\end{lemma}

\begin{theorem}[Outer Relaxation Round]
\lean{Complexity.LogProg.ExpBFS.outer_run}
\label{Complexity.LogProg.ExpBFS.outer_run}
\leanok
\uses{Complexity.LogProg.ExpBFS.A, Complexity.LogProg.ExpBFS.G, Complexity.LogProg.ExpBFS.Mono, Complexity.LogProg.ExpBFS.Mono.refl, Complexity.LogProg.ExpBFS.Mono.trans, Complexity.LogProg.ExpBFS.Out, Complexity.LogProg.ExpBFS.Out.refl, Complexity.LogProg.ExpBFS.Out.trans, Complexity.LogProg.ExpBFS.SBd, Complexity.LogProg.ExpBFS.Site.idx, Complexity.LogProg.ExpBFS.Site.set, Complexity.LogProg.ExpBFS.St, Complexity.LogProg.ExpBFS.St.Bd, Complexity.LogProg.ExpBFS.St.toV_U, Complexity.LogProg.ExpBFS.acExit, Complexity.LogProg.ExpBFS.ac_run, Complexity.LogProg.ExpBFS.cell, Complexity.LogProg.ExpBFS.cf, Complexity.LogProg.ExpBFS.cont, Complexity.LogProg.ExpBFS.ex2_pre, Complexity.LogProg.ExpBFS.inner_run, Complexity.LogProg.ExpBFS.mono_antisymm, Complexity.LogProg.ExpBFS.reach_cf, Complexity.LogProg.ExpBFS.rr, Complexity.LogProg.ExpBFS.step_of, Complexity.LogProg.ExpBFS.stp, Complexity.LogProg.MReach, Complexity.LogProg.MReach.refl, Complexity.LogProg.MReach.trans, Complexity.LogProg.astep}
\difficulty{7}
Let $K_0,B,P$ be natural numbers, let $\mathrm{orc}$ assign a Boolean to each of four
indices and each Boolean string, let $x$ be a Boolean string, and let $\varphi$ be a
Boolean. Let $j$ be a natural number, $S$ a machine state, and
$M:\mathbb{Z}\to\{\mathrm{false},\mathrm{true}\}$ a memory. Suppose $S.K-S.U=j$, $S.U\le
S.K$, $S.V=S.T=S.T_2=S.Z=0$, $S.ch\le1$, and $S$ satisfies the stated register bounds
$\mathrm{SBd}(B,P)$. Then there are a memory $M'$ and a natural number $ch'$ such that
the machine reaches a configuration at the round-exit instruction with $U=S.K$, change
flag $ch'$, and memory $M'$, with the guard $G(x,B,P)$ holding at every configuration
before the endpoint. The memory $M'$ extends the marked positions of $M$ and agrees with
$M$ outside the cells $S.B_s+v$ for $v<S.K$; every predicate on indices below $S.K$ that
is preserved along the step relation and holds at every initially marked index below
$S.K$ also holds at every index below $S.K$ marked in $M'$. If $M'=M$, every step edge
from an initially marked index $a$ with $S.U\le a<S.K$ to an index $b<S.K$ has $b$
marked in $M$, and $ch'=S.ch$; if $M'\ne M$, then $ch'=1$. In either case $ch'\le1$.
\end{theorem}
